\scanpage{120}
% letterhead
COLLÈGE\\
DE\\
FRANCE\\
\rule{1em}{.4pt}\\
\deleted{\textit{Paris, le}}xxxxx Massy le 14.9.1971\\
Cher Monsieur Poncet,\\
Je viens d'avoir connaissance de deux ennuis sérieux avec la machine Hermès 3000 que vous m'aviez livrée en juin: l\deleted{\unsure{a}}e margeur gauche ne fonctionne plus, et l'enroulement du ruban est défectueux, ne permettant d'utiliser qu'une partie du ruban et obligeant de ce fait à des manipulations supplémentaires fort gênantes. Ces ennuis se présentent alors que la machine a encore à peine servi (une trentaine de pages de texte environ), et il est clair que la qualité des machines Hermès 3000 a considérablement baissée depuis une dizaine d'années (j'ai une machine avec laquelle je n'ai eu aucun ennui depuis cette période). C'est pourquoi je me vois dans l'obligation \inserted{d'annuler} la commande d'une autre machine Hermès 3000 que je vous avais faite vers la fin du mois d'août (lettre datée erronnément du 1.6.1971 ...), préférant m'adresser à une marque de machine portative qui ne m'exposera pas à des contretemps continu\deleted{\unsure{le|el}}\unsure{el|le}s\unsure{.}

Au cas où à la réception de la présente lettre, la livraison aurait déjà été faite, je vous serais obligé de la reprendre au Collège de France. Je vous renvoie ci-joints facture et bon de garantie relatifs à cette commande.

Veuillez agréer, cher Monsieur, l'expression de ma considération\\
la plus distinguée\\
% archive sheet number: 119
